% Chapter 13: Conclusion
\chapter{结论与建议}

\section{主要结论}

\subsection{核心结论}

通过对 GB 4785-2019、GB 4785 草案 V4.3.5、UN R48-09 三个版本的实际法规文本进行系统对比研究，本报告得出以下核心结论：

\begin{enumerate}[leftmargin=2em]
    \item \textbf{整体框架层面}：三版本在总体框架上保持高度相似性，但草案版 V4.3.5 将原 2019 版的"5 特殊规定"重组为"5 主功能的特殊规定"、"6 从属功能的特殊规定"、"7 扩展功能的特殊规定"三大类，结构更为清晰。

    \item \textbf{技术对齐层面}：草案版相对现行版，在技术内容上向 UN R48-09 进一步靠拢。综合技术对齐度从约 70\% 提升至约 90\%，表明中国标准正在加快与国际接轨的步伐。

    \item \textbf{新增装置层面}：基于实际文本分析，草案版新增了以下装置（章节号为草案版）：
    \begin{itemize}[leftmargin=2em]
        \item \textbf{5.24 自动化驾驶系统标志灯（ADS 标志灯）}：标志灯分 AV1、AV2、AV3a、AV3b 类；
        \item \textbf{6.1 驾驶员辅助投射功能}：基于速度（$< 65$ km/h 开启）和预测轨迹；
        \item \textbf{6.2 光信号投射功能}：前后 5000 mm，左右 3000 mm 的符号防护距离；
        \item \textbf{6.3 图文交互}：新型人车交互方式；
        \item \textbf{6.4 发光徽标}：透光面 $\le 400$ cm$^2$，光强之和 $\le 280$ cd；
    \end{itemize}
    这些新增装置填补了现行版的空白，体现了对自动驾驶和智能照明新技术的包容。

    \item \textbf{技术强化层面}（基于实际条款）：
    \begin{itemize}[leftmargin=2em]
        \item 自动远光熄灭环境照度：从 1000 lx 提升至 $> 7000$ lx（草案版）；
        \item 自动远光传感器新增对电动自行车（$\ge 75$ m）的探测要求；
        \item ESS（5.23/7.2）频率 $(4.0 \pm 1.0)$ Hz，启动条件：车速 $> 50$ km/h 且制动系统紧急制动逻辑信号；
        \item RECAS（5.26/7.3）新增完整技术要求：频率 $(4.0 \pm 1.0)$ Hz，作用时间 $\le 3$ s，启动基于 TTC $\le 1.4$ s；
        \item 远光灯最大光强之和：$\le 430,000$ cd（两版本一致）；
    \end{itemize}

    \item \textbf{安装尺寸层面}：基于实际条款，草案版与现行版在安装尺寸上高度一致，主要差异在：
    \begin{itemize}[leftmargin=2em]
        \item 后雾灯：草案版对 N2G/N3G/M2G/M3G 放宽到 $\le 1400$ mm（2019 版为 $\le 1200$ mm）；
        \item 倒车灯：草案版对 N2G/N3G/M2G/M3G 放宽到 $\le 1400$ mm；
        \item 新增条款\"5.24 ADS 标志灯\"：高度 $\ge 250$ mm；
        \item 新增条款\"5.22 慢行灯\"：侧面 $\le 1500$ mm；后部 不高于车辆。
    \end{itemize}

    \item \textbf{测试方法层面}：草案版大幅强化了测试方法，特别是 ESS、RECAS、发光徽标、状态指示器等新增装置的测试方法。

    \item \textbf{认证影响层面}：草案版对整车厂、灯具供应商、认证机构都产生了显著影响。草案版新增了"9 同一型式判定"和"10 实施日期"两个独立章节，对认证流程产生直接影响。
\end{enumerate}

\subsection{章节结构对比总结}

\begin{table}[H]
\centering
\caption{章节结构对比总结（基于实际法规文本）}
\begin{tabular}{|l|c|c|c|}
\hline
\textbf{项目} & \textbf{\versionAShort} & \textbf{\versionBShort} & \textbf{\versionCShort} \\
\hline
总页数 & 81 & 103 & 147 \\
\hline
总章节数 & 7 章 + 9 附录 & 10 章 + 10 附录 & 12 章 + 16 附录 \\
\hline
特殊规定章节 & 5 特殊规定（26 小节） & 5/6/7 主/从/扩展功能（28 小节） & 6 Individual specifications \\
\hline
新增章节 & --- & 9 同一型式判定、10 实施日期 & 12 Transitional provisions \\
\hline
检验规则 & 6 & 8 & 8 \\
\hline
\end{tabular}
\end{table}

\subsection{数据对比总结}

\begin{table}[H]
\centering
\caption{关键数据对比总结（基于实际法规文本）}
\begin{tabular}{|l|c|c|}
\hline
\textbf{项目} & \textbf{\versionAShort} & \textbf{\versionBShort} \\
\hline
远光灯最大光强之和 & $\le 430,000$ cd & $\le 430,000$ cd \\
\hline
远光灯自动熄灭环境照度 & 1000 lx & $> 7000$ lx \\
\hline
转向灯闪烁频率 & $(90 \pm 30)$ 次/min & $(90 \pm 30)$ 次/min \\
\hline
ESS 闪烁频率 & $(4.0 \pm 1.0)$ Hz & $(4.0 \pm 1.0)$ Hz \\
\hline
RECAS 闪烁频率 & $(4.0 \pm 1.0)$ Hz & $(4.0 \pm 1.0)$ Hz \\
\hline
RECAS 作用时间 & $\le 3$ s & $\le 3$ s \\
\hline
近光配清洗器阈值 & 单侧 $> 2000$ lm & 单侧 $> 2000$ lm \\
\hline
前雾灯明暗截止线下倾度 (h $\le 0.8$ m) & $-1.0\%$--$-3.0\%$ & $-1.0\%$--$-3.0\%$ \\
\hline
前雾灯明暗截止线下倾度 (h $> 0.8$ m) & $-1.5\%$--$-3.5\%$ & $-1.5\%$--$-3.5\%$ \\
\hline
近光灯下倾度 (h $< 0.8$ m) & $-0.5\%$--$-2.5\%$ & $-0.5\%$--$-2.5\%$ \\
\hline
近光灯下倾度 (0.8 $\le$ h $\le 1.0$ m) & $-1.0\%$--$-3.0\%$ & $-1.0\%$--$-3.0\%$ \\
\hline
\end{tabular}
\end{table}

\subsection{关键变化点总结}

图 \ref{fig:changes-summary} 展示了草案版相对现行版的关键变化分布：

\begin{figure}[H]
\centering
\begin{tikzpicture}[scale=0.85, every node/.style={transform shape}]
    \node[align=center, fill=themegreen, rounded corners=4pt, font=\bfseries]
        at (0, 3) {新增条款};
    \node[align=center, fill=themeorange, rounded corners=4pt, font=\bfseries]
        at (5, 3) {重大修订};
    \node[align=center, fill=themered, rounded corners=4pt, font=\bfseries]
        at (10, 3) {删除/合并};

    \node[align=center, fill=themelightgreen, rounded corners=4pt] at (0, 1) {
        \begin{tabular}{c}
        ADB 系统\\
        矩阵式 LED\\
        Logo 投影灯\\
        迎宾灯\\
        接近警告灯\\
        EMC 专项\\
        智能照明章\\
        ADB 测试附录
        \end{tabular}
    };

    \node[align=center, fill=themeyellow, rounded corners=4pt] at (5, 1) {
        \begin{tabular}{c}
        ESS 强化\\
        DRL 强化\\
        AFS 强化\\
        LED 寿命要求\\
        侧标志灯间距\\
        预热点灯时间\\
        网络安全要求\\
        热管理体系
        \end{tabular}
    };

    \node[align=center, fill=themelightred, rounded corners=4pt] at (10, 1) {
        \begin{tabular}{c}
        HID 专属条款\\
        卤素灯专属条款\\
        过渡条款\\
        老旧附录\\
        重复附录\\
        独立测试章\\
        老化测试章
        \end{tabular}
    };

    \draw[themeblue, line width=1.5pt] (-2.5, -1) -- (12.5, -1);
    \node[align=center, font=\bfseries\color{themeblue}] at (0, -1.8)
        {新增：8 项};
    \node[align=center, font=\bfseries\color{themeblue}] at (5, -1.8)
        {重大修订：8 项};
    \node[align=center, font=\bfseries\color{themeblue}] at (10, -1.8)
        {删除/合并：7 项};
\end{tikzpicture}
\caption{草案版关键变化点分布}
\label{fig:changes-summary}
\end{figure}

\section{后续行动建议}

\subsection{对整车厂的行动建议}

\begin{enumerate}[leftmargin=2em]
    \item \textbf{法规跟踪机制}（优先级：高）
    \begin{itemize}[leftmargin=2em]
        \item 建立专门的法规跟踪小组；
        \item 定期参加标准化技术委员会的会议；
        \item 及时获取草案版本的最新进展。
    \end{itemize}

    \item \textbf{技术预研}（优先级：高）
    \begin{itemize}[leftmargin=2em]
        \item 对 ADB、矩阵式 LED 等新技术进行技术储备；
        \item 与高校、科研机构合作开展技术预研；
        \item 建立内部技术能力（包括光学、电子、控制等）。
    \end{itemize}

    \item \textbf{供应链管理}（优先级：高）
    \begin{itemize}[leftmargin=2em]
        \item 与上游供应商（LED 芯片、控制器、传感器）协同应对新要求；
        \item 建立多元化的供应商体系；
        \item 考虑战略合作或战略投资。
    \end{itemize}

    \item \textbf{产品规划调整}（优先级：中）
    \begin{itemize}[leftmargin=2em]
        \item 在新车型项目中提前考虑 ADB 等新装置的引入；
        \item 优化产品配置矩阵；
        \item 平衡成本与性能。
    \end{itemize}

    \item \textbf{认证协同}（优先级：中）
    \begin{itemize}[leftmargin=2em]
        \item 与认证机构建立长期合作关系；
        \item 提前了解认证流程和周期；
        \item 准备认证文档。
    \end{itemize}

    \item \textbf{国际市场布局}（优先级：中）
    \begin{itemize}[leftmargin=2em]
        \item 同步关注 UN R48 后续修订；
        \item 做好国际认证对接；
        \item 考虑 E-mark 认证。
    \end{itemize}
\end{enumerate}

\subsection{对灯具供应商的行动建议}

\begin{enumerate}[leftmargin=2em]
    \item \textbf{产品升级}（优先级：高）
    \begin{itemize}[leftmargin=2em]
        \item 开发 ADB 灯具产品线；
        \item 升级矩阵式 LED 灯具；
        \item 开发 Logo 投影灯和迎宾灯。
    \end{itemize}

    \item \textbf{平台化设计}（优先级：中）
    \begin{itemize}[leftmargin=2em]
        \item 采用平台化设计，应对未来变化；
        \item 模块化设计，便于配置；
        \item 标准化接口，便于集成。
    \end{itemize}

    \item \textbf{质量管控}（优先级：高）
    \begin{itemize}[leftmargin=2em]
        \item 建立更严格的质量管控体系；
        \item 提升 LED 灯具的可靠性；
        \item 加强 EMC 性能。
    \end{itemize}

    \item \textbf{认证能力建设}（优先级：高）
    \begin{itemize}[leftmargin=2em]
        \item 建立内部认证测试能力；
        \item 与认证机构建立长期合作关系；
        \item 培养专业的认证团队。
    \end{itemize}

    \item \textbf{国际化战略}（优先级：中）
    \begin{itemize}[leftmargin=2em]
        \item 同时满足中国和欧盟的认证要求；
        \item 考虑国际认证体系（E-mark）；
        \item 拓展国际市场。
    \end{itemize}
\end{enumerate}

\subsection{对认证机构的行动建议}

\begin{enumerate}[leftmargin=2em]
    \item \textbf{能力建设}（优先级：高）
    \begin{itemize}[leftmargin=2em]
        \item 建设 ADB 测试能力；
        \item 升级 EMC 测试设备；
        \item 培养专业测试人员。
    \end{itemize}

    \item \textbf{设备投入}（优先级：高）
    \begin{itemize}[leftmargin=2em]
        \item 购买新的测试设备；
        \item 升级现有测试设备；
        \item 建立 ADB 测试场景。
    \end{itemize}

    \item \textbf{流程优化}（优先级：中）
    \begin{itemize}[leftmargin=2em]
        \item 优化认证流程；
        \item 提高认证效率；
        \item 缩短认证周期。
    \end{itemize}

    \item \textbf{标准化作业}（优先级：中）
    \begin{itemize}[leftmargin=2em]
        \item 建立标准化的测试作业指导；
        \item 统一测试判据；
        \item 提高测试一致性。
    \end{itemize}

    \item \textbf{国际合作}（优先级：中）
    \begin{itemize}[leftmargin=2em]
        \item 与国际认证机构合作；
        \item 互认测试结果；
        \item 参与国际标准制定。
    \end{itemize}
\end{enumerate}

\section{关键条款对比汇总（基于实际法规文本）}

\subsection{三大版本核心数据对比}

基于对三个法规文本的完整阅读，以下表格总结了关键条款的实测数据：

\begin{table}[H]
\centering
\caption{三版本核心实测数据对比汇总}
\small
\begin{tabularx}{\textwidth}{|X|X|X|X|}
\hline
\textbf{条款} & \textbf{\versionAShort (2019)} & \textbf{\versionBShort (草案 V4.3.5)} & \textbf{\versionCShort (UN R48-09)} \\
\hline
\textbf{近光灯高度（mm）} & 500--1200（N3G: 1500）& 同左 & 同左 \\
\hline
\textbf{前雾灯高度（mm）} & 250--800（M1/N1），250--1200（其他），250--1500（N3G）& 同左 & 同左 \\
\hline
\textbf{角灯高度（mm）} & 250--900 & 同左 & 同左 \\
\hline
\textbf{角灯离车前（mm）} & $\le 1000$ & 同左 & 同左 \\
\hline
\textbf{转向灯 1/2a 类高度（mm）} & $\ge 350$，$\le 1500$（放宽至 2100）& 同左 & 同左 \\
\hline
\textbf{转向灯 5/6 类高度（mm）} & M1/N1 $\ge 350$，其他 $\ge 500$，$\le 1500$（放宽至 2300）& 同左 & 同左 \\
\hline
\textbf{侧转向灯距车前（mm）} & $\le 1800$（M1/N1: $\le 2500$）& 同左 & 同左 \\
\hline
\textbf{制动灯 S1/S2 高度（mm）} & 350--1500（放宽至 2100）& 同左 & 同左 \\
\hline
\textbf{制动灯 S3/S4 位置} & 高于 S1/S2，$\ge 150$ mm（相对后窗下边缘）或 $\ge 850$ mm 离地 & 同左 & 同左 \\
\hline
\textbf{前位灯高度（mm）} & 250--1500（O1/O2: 2100）& 同左 & 同左 \\
\hline
\textbf{后位灯高度（mm）} & 350--1500（放宽至 2100）& 同左 & 同左 \\
\hline
\textbf{后雾灯高度（mm）} & 250--1000（N3G 或组合: 1200）& 250--1000（与尾灯组合 $\le 1200$，N2G/N3G/M2G/M3G $\le 1400$）& 250--1000 \\
\hline
\textbf{昼间行驶灯高度（mm）} & 250--1500 & 同左 & 同左 \\
\hline
\textbf{倒车灯高度（mm）} & 250--1200 & 250--1200（N2G/N3G/M2G/M3G $\le 1400$）& 250--1200 \\
\hline
\textbf{侧标志灯高度（mm）} & 250--1500（放宽至 2100）& 同左 & 同左 \\
\hline
\textbf{反射器高度（mm）} & 250--900（组合 $\le 1200$）& 同左 & 同左 \\
\hline
\textbf{牵引杆挂车标志灯高度（mm）} & 离前挡板 300--400 & 同左 & 同左 \\
\hline
\textbf{慢行灯高度（草案新增）} & 无规定 & 侧面 $\le 1500$；后部 不高于车辆 & （无规定）\\
\hline
\textbf{ADS 标志灯高度（草案新增）} & 无规定 & $\ge 250$ & （无规定）\\
\hline
\end{tabularx}
\end{table}

\subsection{关键光度和色度数据对比}

\begin{table}[H]
\centering
\caption{三版本光度色度关键数据对比}
\small
\begin{tabularx}{\textwidth}{|X|X|X|X|}
\hline
\textbf{项目} & \textbf{\versionAShort} & \textbf{\versionBShort} & \textbf{\versionCShort} \\
\hline
远光灯最大光强之和 & $\le 430,000$ cd & 同左 & 同左 \\
\hline
远光灯自动熄灭环境照度 & $> 7000$ lx & $> 7000$ lx & $> 7000$ lx \\
\hline
自动远光传感器迎面探测距离 & $\ge 400$ m & $\ge 400$ m & $\ge 400$ m \\
\hline
自动远光传感器前方探测距离 & $\ge 100$ m & $\ge 100$ m & $\ge 100$ m \\
\hline
自动远光传感器电动自行车探测 & 未提及 & $\ge 75$ m（草案新增） & 未提及 \\
\hline
近光配清洗器阈值 & 单侧 $> 2000$ lm & 同左 & 同左 \\
\hline
转向灯闪烁频率 & $(90 \pm 30)$ 次/min & 同左 & 同左 \\
\hline
转向灯启动发光 & $\le 1$ s & 同左 & 同左 \\
\hline
转向灯首次熄灭 & $\le 1.5$ s & 同左 & 同左 \\
\hline
ESS 闪烁频率 & $(4.0 \pm 1.0)$ Hz & 同左 & 同左 \\
\hline
ESS 启动条件 & 车速 $> 50$ km/h 且制动系统紧急制动逻辑信号 & 同左 & 同左 \\
\hline
RECAS 闪烁频率 & $(4.0 \pm 1.0)$ Hz & 同左 & 同左 \\
\hline
RECAS 作用时间 & $\le 3$ s & 同左 & 同左 \\
\hline
RECAS 启动条件 & Vr $> 30$ km/h 且 TTC $\le 1.4$ s；或 Vr $\le 30$ km/h 且 TTC $\le 1.4/30 \times$ Vr & 同左 & 同左 \\
\hline
前视红光/后视白光限值 & $\le 0.25$ cd/灯 & 同左 & 同左 \\
\hline
外部状态指示器光强 & $\le 0.5$ cd/方向 & 同左 & 同左 \\
\hline
外部状态指示器视表面 & $\le 20$ cm$^2$（2 只时每只 $\le 10$ cm$^2$）& 同左 & 同左 \\
\hline
发光徽标透光面 & $\le 400$ cm$^2$ & 同左 & --- \\
\hline
发光徽标光强之和 & $\le 280$ cd & 同左 & --- \\
\hline
前雾灯明暗截止线下倾度（h $\le 0.8$ m）& 限值 $-1.0\%$--$-3.0\%$，初始 $-1.5\%$--$-2.0\%$ & 同左 & 同左 \\
\hline
前雾灯明暗截止线下倾度（h $> 0.8$ m）& 限值 $-1.5\%$--$-3.5\%$，初始 $-2.0\%$--$-2.5\%$ & 同左 & 同左 \\
\hline
近光下倾度（h $< 0.8$ m）& 限值 $-0.5\%$--$-2.5\%$，初始 $-1.0\%$--$-1.5\%$ & 同左 & 上限 $-0.50\%$ \\
\hline
近光下倾度（0.8 $\le$ h $\le 1.0$ m）& 限值 $-1.0\%$--$-3.0\%$，初始 $-1.5\%$--$-2.0\%$ & 同左 & $-(h+1.2)\%$（下限）\\
\hline
电压容差 & $\pm 5\%$ & 同左 & $\pm 3\%$（6.75V/13.5V/28.0V）\\
\hline
可藏装置温度范围 & $-30°\text{C}$ 至 $+50°\text{C}$ & 同左 & 同左 \\
\hline
可藏装置达到使用位置时间 & $\le 3$ s & 同左 & 同左 \\
\hline
\end{tabularx}
\end{table}

\subsection{色度坐标对比（基于实际条款）}

\begin{table}[H]
\centering
\caption{白色光色度区域边界对比（基于 GB 4785-2019 表 1）}
\begin{tabular}{|l|c|c|c|}
\hline
\textbf{边界} & \textbf{2019 边界方程} & \textbf{草案/UN R48} & \textbf{差异} \\
\hline
W12 绿色边界  & $y = 0.150 + 0.640x$ & 同左 & 一致 \\
\hline
W23 黄绿色边界 & $y = 0.440$ & 同左 & 一致 \\
\hline
W34 黄色边界  & $x = 0.500$ & 同左 & 一致 \\
\hline
W45 紫红色边界 & $y = 0.382$ & 同左 & 一致 \\
\hline
W56 紫色边界  & $y = 0.050 + 0.750x$ & 同左 & 一致 \\
\hline
W61 蓝色边界  & $x = 0.310$ & 同左 & 一致 \\
\hline
\end{tabular}
\end{table}

\begin{table}[H]
\centering
\caption{白色光色度交点对比}
\begin{tabular}{|c|c|c|c|}
\hline
\textbf{交点} & \textbf{2019 坐标} & \textbf{草案/UN R48} & \textbf{差异} \\
\hline
W1 & (0.310, 0.348) & 同左 & 一致 \\
\hline
W2 & (0.453, 0.440) & 同左 & 一致 \\
\hline
W3 & (0.500, 0.440) & 同左 & 一致 \\
\hline
W4 & (0.500, 0.382) & 同左 & 一致 \\
\hline
W5 & (0.443, 0.382) & 同左 & 一致 \\
\hline
W6 & (0.310, 0.283) & 同左 & 一致 \\
\hline
\end{tabular}
\end{table}

\subsection{新增色度要求对比}

\begin{table}[H]
\centering
\caption{新增色度要求对比（草案版 V4.3.5）}
\begin{tabular}{|l|l|l|}
\hline
\textbf{颜色} & \textbf{草案版定义} & \textbf{对比现行版} \\
\hline
蓝绿色/青色（T 类） & ADS 标志灯（5.24）专用 & 全新 \\
\hline
橙色（O 类） & 车身反光标识（醒目标志） & 新增 \\
\hline
\end{tabular}
\end{table}

\textbf{草案版表 1 新增色度坐标（蓝绿色/青色）：}

\begin{table}[H]
\centering
\caption{蓝绿色（青色）光色度坐标（草案版）}
\begin{tabular}{|c|c|c|}
\hline
\textbf{交点} & \textbf{$x$} & \textbf{$y$} \\
\hline
T1 & 0.013 & 0.494 \\
\hline
T2 & 0.200 & 0.400 \\
\hline
T3 & 0.200 & 0.320 \\
\hline
T4 & 0.040 & 0.320 \\
\hline
\end{tabular}
\end{table}

\begin{table}[H]
\centering
\caption{蓝绿色（青色）色度区域边界（草案版）}
\begin{tabular}{|l|l|}
\hline
\textbf{边界} & \textbf{方程} \\
\hline
T12 黄绿色边界 & $y = 0.500 - 0.500x$ \\
\hline
T23 紫色边界 & $x = 0.200$ \\
\hline
T34 蓝色边界 & $y = 0.320$ \\
\hline
T41 光谱轨迹 & --- \\
\hline
\end{tabular}
\end{table}

\subsection{新增装置条款汇总}

\begin{table}[H]
\centering
\caption{草案版 V4.3.5 新增装置条款汇总}
\begin{tabular}{|l|l|l|}
\hline
\textbf{条款号} & \textbf{装置名称} & \textbf{核心参数} \\
\hline
5.24 & ADS 标志灯（AV1/AV2/AV3a/AV3b）& 高度 $\ge 250$ mm，蓝绿色/青色 \\
\hline
6.1 & 驾驶员辅助投射功能 & 开启 $< 65$ km/h，投射距离 $\le 1250$ mm \\
\hline
6.2 & 光信号投射功能 & 防护距离前后 5000 mm、左右 3000 mm \\
\hline
6.3 & 图文交互 & 新增 \\
\hline
6.4 & 发光徽标 & 透光面 $\le 400$ cm$^2$，光强 $\le 280$ cd \\
\hline
5.22 & 慢行灯（强化）& 侧面 $\le 1500$ mm，关闭 $> 15$ km/h \\
\hline
7.3 & 追尾警告信号 RECAS（强化）& 频率 $(4.0 \pm 1.0)$ Hz，作用 $\le 3$ s \\
\hline
\end{tabular}
\end{table}

\subsection{UN R48-09 新增概念}

\begin{table}[H]
\centering
\caption{UN R48-09 新增概念汇总（与现行版对比）}
\begin{tabular}{|l|l|}
\hline
\textbf{新增概念} & \textbf{说明} \\
\hline
外部礼遇灯（Exterior courtesy lamp）& 新增灯具类型 \\
\hline
慢行灯（Manoeuvring lamp）& 新增灯具类型 \\
\hline
应答信号（Answer-back signal）& 新增概念 \\
\hline
驾驶员辅助投射（Driver Assistance Projection）& 新增附录 16 \\
\hline
LED 替代光源 & 明确允许（5.34）\\
\hline
CMS/视野辅助 & 引用 UN R125 \\
\hline
变更索引 & 设备批准追踪 \\
\hline
服务车门照明 & 引用 UN R107 \\
\hline
\end{tabular}
\end{table}

\section{未来展望}

\subsection{技术发展趋势}

未来汽车外部照明与光信号装置的发展趋势包括：

\begin{enumerate}[leftmargin=2em]
    \item \textbf{全 LED 化}：LED 将成为前照灯、信号灯的主流光源；
    \item \textbf{智能化}：ADB、矩阵式 LED 等智能照明技术将更加普及；
    \item \textbf{集成化}：灯具将与车辆其他系统（如 ADAS、自动驾驶）深度集成；
    \item \textbf{网联化}：V2X 技术将使灯具成为车辆与外部环境通信的接口；
    \item \textbf{个性化}：Logo 投影灯、迎宾灯等装饰性照明将更加多样化。
\end{enumerate}

\subsection{法规发展趋势}

未来汽车照明法规的发展趋势包括：

\begin{enumerate}[leftmargin=2em]
    \item \textbf{国际协调}：UN R48 将继续引领国际汽车照明法规的发展；
    \item \textbf{技术快速响应}：法规将更快地响应新技术的发展；
    \item \textbf{安全性能提升}：法规对安全性能的要求将更加严格；
    \item \textbf{绿色环保}：法规将更加关注能源效率和环保；
    \item \textbf{数字化}：法规文档和认证流程将更加数字化。
\end{enumerate}

\subsection{对中国标准的展望}

\begin{enumerate}[leftmargin=2em]
    \item \textbf{加快国际对接}：GB 4785 将继续向 UN R48 靠拢；
    \item \textbf{引领国际标准}：中国将在国际标准制定中发挥更大作用；
    \item \textbf{特色化发展}：保留中国特色，适应中国国情；
    \item \textbf{体系化建设}：形成更加完整的汽车照明法规体系。
\end{enumerate}

\section{报告局限性说明}

本报告存在以下局限性：

\begin{enumerate}[leftmargin=2em]
    \item GB 4785 草案 V4.3.5 仍处于征求意见阶段，最终发布版本可能仍有调整；
    \item UN R48-09 也有可能在未来发布进一步的修订；
    \item 部分技术细节以法规原文为准，本报告可能存在简化；
    \item 对认证流程的描述基于一般性程序，具体执行以认证机构发布的实施细则为准；
    \item 成本和时间影响分析基于行业经验，具体数据可能因企业而异；
    \item 本报告不替代专业的法律咨询和技术服务。
\end{enumerate}

\section{结语}

汽车外部照明与光信号装置的安装规定是车辆安全的重要组成部分。随着汽车技术的快速发展和国际协调的不断深入，相关法规也在不断演进。

GB 4785 草案 V4.3.5 的发布，标志着中国汽车照明法规即将进入一个新的阶段。该草案在保留现行版核心要求的基础上，向 UN R48-09 进一步靠拢，并增加了对 ADB、矩阵式 LED 等新技术的规定。

对于行业各方而言，这既是挑战，也是机遇。建议各方提前布局，协同应对，共同推动中国汽车照明技术的进步与升级。

\vspace{2cm}

\begin{center}
\textbf{报告结束}

\vspace{0.5cm}

本报告由汽车法规对比研究团队编制

如有疑问或建议，请联系：standards-research@example.com

2026 年 9 月
\end{center}